\section{Discussion}

The central result is a separation between favorable point estimates and
stable external evidence. Each candidate produced a positive mean paired
Pearson difference, and each passed the adjusted one-sided randomization
threshold. Two candidates improved on every observed seed. Yet the
hierarchical intervals remained compatible with zero once both seed and subject
variation were propagated. Under the predeclared conjunction, no tested complex
head therefore established a stable gain.

This finding changes how the head comparison should be interpreted. The
earlier-layer linear result suggests that age-relevant signal is not confined to
the final REVE layer. The rich-statistics head suggests that token dispersion
and magnitude can be useful to a downstream predictor. The multi-query result
shows that additional adaptive aggregation can improve error metrics without
delivering a reliable Pearson advantage. None of these observations warrants
the stronger claim that the frozen representation itself is better under a more
expressive readout.

The experiment also illustrates why probe capacity should be treated as part of
the scientific intervention. A nonlinear head can recover target information
that is present but not linearly accessible, but it can also exploit
development-set regularities that do not transfer. Matched seeds, identical
cached representations, common optimization, and paired external subjects
reduce avoidable comparison noise; they do not eliminate cohort uncertainty.

Finally, the negative confirmatory outcome is informative rather than a failed
experiment. It narrows the supported claim to the tested REVE checkpoint,
preprocessing, HBN development split, MIPDB cohort, head families, and seeds. It
motivates future work on larger independent cohorts, cross-cohort replication,
sample-efficiency curves, frequency-specific perturbations, and explicitly
regularized probes. Such extensions should be specified as new studies rather
than retrofitted to the sealed analysis.
